\documentclass[../../main.tex]{subfiles}
\begin{document}

{\bf [解答]}

\begin{figure}[htb]
  \centering
  \begin{tikzpicture}[x=1.45cm, y=1.05cm, >=stealth, every node/.style={font=\small}]
    % Illustrative alpha = -0.65, beta = 1.35 on y = x^2
    \pgfmathsetmacro{\xa}{-0.65}
    \pgfmathsetmacro{\ya}{\xa*\xa}
    \pgfmathsetmacro{\xb}{1.35}
    \pgfmathsetmacro{\yb}{\xb*\xb}
    \pgfmathsetmacro{\xm}{0.5*(\xa+\xb)}
    \pgfmathsetmacro{\ym}{0.5*(\ya+\yb)}
    \coordinate (O) at (0, 0);
    \coordinate (A) at (\xa, \ya);
    \coordinate (B) at (\xb, \yb);
    \coordinate (M) at (\xm, \ym);

    \draw[->, thick, gray!80!black] (-1.65, 0) -- (1.85, 0) node[right] {$x$};
    \draw[->, thick, gray!80!black] (0, -0.35) -- (0, 2.55) node[above] {$y$};

    % 放物線 y = x^2
    \draw[line width=1.4pt, blue!85!black, domain=-1.5:1.55, samples=80, smooth]
    plot (\x, {\x*\x}) node[above right=1pt, font=\footnotesize] {$y=x^2$};

    % 線分 AB (長さ l)
    \draw[line width=1.3pt, teal!85!black] (A) -- (B);
    \node[above left=2pt, font=\footnotesize, teal!85!black] at ($(A)!0.32!(B)$) {$l$};

    % 投影線
    \draw[dashed, gray!65] (A) -- (\xa, 0) node[below=2pt, font=\footnotesize, black] {$\alpha$};
    \draw[dashed, gray!65] (B) -- (\xb, 0) node[below=2pt, font=\footnotesize, black] {$\beta$};
    \draw[dashed, red!75!black] (M) -- (\xm, 0) node[below=2pt, font=\footnotesize, red!85!black] {$X$};

    % 点とラベル
    \node[circle, fill=teal!85!black, inner sep=1.3pt] at (A) {};
    \node[left=4pt, font=\footnotesize, teal!85!black] at (A) {$A(\alpha,\alpha^2)$};
    \node[circle, fill=teal!85!black, inner sep=1.3pt] at (B) {};
    \node[right=3pt, font=\footnotesize, teal!85!black] at (B) {$B(\beta,\beta^2)$};
    \node[circle, fill=red!85!black, inner sep=1.4pt] at (M) {};
    \node[above=3pt, font=\footnotesize, red!85!black, fill=white, inner sep=1pt] at (0.48, 1.18) {$M(X,Y)$};
    \node[below left=1pt, font=\footnotesize] at (O) {$O$};
  \end{tikzpicture}
  \caption{放物線 $y=x^2$ 上に両端をのせて動く長さ $l$ の線分 $AB$ とその中点 $M(X,Y)$}
  \label{fig:1974_2_overview}
\end{figure}

線分の両端を$A(\alpha,\alpha^2)$，$B(\beta,\beta^2)$とする．ただし
\begin{align}
  \alpha<\beta\label{eq:1}
\end{align}
とする．$|AB|=l$だから
\begin{align}
  |AB|=l \Longleftrightarrow |AB|^2=l^2  \nonumber\\
  (\beta-\alpha)^2+(\beta^2-\alpha^2)^2=l^2  \nonumber\\
  (\beta-\alpha)^2\left(1+(\alpha+\beta)^2\right)=l^2 \label{eq:2}
\end{align}
となる．
ここで$a=\alpha+\beta$，$b=\beta-\alpha$とおいて，$a$，$b$の存在条件を調べる．
まずは\cref{eq:1,eq:2}に代入して
\begin{align}
  \begin{cases}
    b > 0 \\
    b^2(1+a^2)=l^2
  \end{cases} \label{eq:3}
\end{align}
である．

次に$\alpha\beta=\dfrac{a^2-b^2}{4}$であるから，$\alpha$，$\beta$は$x$の$2$次方程式$x^2-ax+\dfrac{a^2-b^2}{4}=0$の異なる$2$実解であるから，判別式$D$として
\begin{align}
  &D>0
  \Longleftrightarrow
  a^2-4\dfrac{a^2-b^2}{4}>0
  \Longleftrightarrow
  b^2>0
\end{align}
である．これは\cref{eq:3}より常に成立するから，$a,b$の条件式は\cref{eq:3}である．

このもとで$M(X,Y)$として$Y$が最小となる場合を考えれば良い($\because Y\ge0$)．
それぞれの座標は
\begin{align}
  \begin{cases}
    X=\dfrac{\alpha+\beta}{2}=\dfrac{a}{2}  \\
    Y=\dfrac{\alpha^2+\beta^2}{2}=\dfrac{a^2+b^2}{4}
  \end{cases} \label{eq:4}
\end{align}
である．
\cref{eq:3}から$b^2=\dfrac{l^2}{1+a^2}$を代入して
\begin{align}
  Y&=\dfrac{a^2+b^2}{4} \\
  &=\dfrac{1}{4}\left(a^2+\dfrac{l^2}{1+a^2}\right)  \\
  &=\dfrac{1}{4}\left(1+a^2+\dfrac{l^2}{1+a^2}-1\right)  \\
  &\ge \dfrac{1}{4}\left(2\sqrt{l^2}-1\right) \quad (\because \text{AM-GM})\\
  &=\dfrac{1}{4}\left(2l-1\right) \quad (\because l>0)
\end{align}
である．等号成立は
\begin{align}
  1+a^2=\dfrac{l^2}{1+a^2} \Longleftrightarrow a = \pm \sqrt{l-1} \quad (\because l \ge 1)
\end{align}
の時である．この時\cref{eq:4}より
\begin{align}
  X=\dfrac{\pm\sqrt{l-1}}{2}
\end{align}
だから，求める座標は
\begin{align}
  \left(\dfrac{\pm\sqrt{l-1}}{2},\dfrac{1}{4}\left(2l-1\right)\right)
\end{align}
である．

\newpage
\end{document}
